\section{Macroeconomics in the Open Economy}



\begin{frame}{The Open Economy}
    \showmytikz[1,2]{../shared/tikz/tikz-flowmodel-open}
    \mycomment{
    We understand the economy of a country to consist chiefly of labour markets, the markets for goods and services, and capital markets.

    The state influences in various ways. It taxes its own citizens, regulates activities, provides currency, and purchases from them.
    }
    
\end{frame}

\begin{frame}{The Open Economy}
    \showmytikz[1,2,3]{../shared/tikz/tikz-flowmodel-open}
    \mycomment{
    But the economy is not alone. There are foreign countries, and there is indeed a world market, where there are also markets for labour, goods and services, and capital. 
    
    Because goods can easily be re-routed if there is arbitrage, in effect there is a highly efficient world market. 
    }
    
\end{frame}

\begin{frame}{The Open Economy}
    \showmytikz[2,4]{../shared/tikz/tikz-flowmodel-open}
    \mycomment{
    Unless there are barriers, and abstracting from transportation costs, the markets for goods and services merge. 

    The markets for labour may or may not merge; but labour is clearly less mobile than goods and capital, due to administrative and language barriers, but also due to family bounds.

    There is, however, a connection here too: firms may decide to produce abroad, investing abroad (foreign direct investment, FDI), and (re-)import.

    Capital markets were not always free, but there are no more capital restrictions at present. Monetary flows are also not taxed. The crossing of goods (and services) across borders might be taxed using tariffs or quotas.
    }
\end{frame}


\begin{frame}{Exchange Rates: Purchasing Power Parity}
    \mycomment{
        In the long run, and when trade flows freely, one would expect the exchange rate to reflect the purchasing power of the currencies. After exchange, I can buy as much in the foreign country as I can buy at home. This can be used to make "real" comparisons. 
    }
    \mytext{
        Exchange rate assures the same basket of goods can be bought here or there.

    (`Big Mac Index')

    \begin{itemize}
        \item often used to `correct' exchange rates 
        \item[] "in USD, real terms"
        \item often used to correct value comparisons over time 
        \item[] "in 2020 USD, real terms"
    \end{itemize}
    }
\end{frame}

\begin{frame}{Exchange Rates: Purchasing Power Parity}
    \mytext{
      Price of haircut in Switzerland 80 SFR $\sim 90 $ \$

    Price of haircut in India: 300 INR (?) $\sim 3.5$ \$

    But: price of notebooks etc. similar

    \begin{itemize}
        \item tradable vs.~non-tradable goods
        \item commodities vs.~manufactured goods
        \item labour productivity / wages
        % \item imbalances / artificial exchange rates
        % \item financial account dominates current account?
    \end{itemize}  
    }
    \mycomment{
        But purchasing power parity (PPP) is always relative to a bundle of goods. If we can trade across borders, then there are world market prices for everything that we trade. There are, however, non-tradable goods.
    }
\end{frame}

\begin{frame}{Exchange Rates: Interest Rate Parity}
    \mycomment{
        Some argue the exchange rate be in reality not dominated by the current account or trade balance. There are also financial flows, and there are larger financial flows than there are flows in goods.

        Then it might indeed be a question of investment. As long as I get more money for my money abroad than at home, I will try to exchange my local currency for the foreign currency and invest abroad.
    }
\mytext{
    Investment in both countries yield is equal
    \begin{equation}
        (1+i_\$) = (1+w) \cdot (1+i_\text{\euro} )
    \end{equation}
    with $i$ interest rate,\\ $w$ expected change of exchange rate;

    \medskip

    Fisher equation:
    \begin{equation}
        (1+i) = (1+r) \cdot  (1+\pi)
    \end{equation}
    with $r$ real exchange rate and $\pi$\ inflation
}
    
\end{frame}

\begin{frame}{Mundell-Fleming Model: Basics}
    \showmytikz{../shared/tikz/tikz-mundellfleming}
    \mycomment{
    The Mundell-Fleming model enlarges the IS-LM model with the introduction of the $ZZ$-curve.
    
    The $ZZ$ curve is also called BoP-Curve. It consists of all combinations of interest and income that balance the Balance of Payments. 

    If capital is perfectly mobile (no capital controls), then the interest rate is given by the international level of interest rates. In the long run, domestic $IS$ and $LM$ curves must adapt.
    
    In this situation, the fiscal and/or monetary policy may or may not be effective, depending on whether the exchange rate is fixed or not, as we will show.

    Monetary policy is now restricted to open market policy, since the central bank can no longer set the interest rate.

    A tariff has the same effect on $IS$ as government spending by rising net exports $\Delta NX$. Thus, all findings relating to fiscal policy apply to tariffs.  
    }
\end{frame}

\begin{frame}{Mundell-Fleming Model: Basics}
    \showmytikz[11]{../shared/tikz/tikz-mundellfleming}
    \mycomment{
        If there are capital controls, the $ZZ$-curve slopes upward. Interest rates will have less and less impact on the Balance of Payments. 

        Capital controls were quite common until as late as the 1980ies, and still exist, in particular in China and Russia. We will not discuss this situation here any further.

        Note that at any point above the $ZZ$-curve, capital flows into the economy, as local interest rates are above the global interest rate. Below, capital flows out. Capital inflows impact the $LM$-curve, pushing money supply up and shifting $LM$ to the right. This is the mechanism that restores the equilibrium.
    }
\end{frame}


\begin{frame}{IS-LM-ZZ Analysis}
    \mycomment{
    We shall now use the Mundell-Fleming model to look at the effectiveness of
    \begin{itemize}
        \item  monetary stimulus through open market activities, i.e. the central bank selling  home currency. (The central bank cannot set the interest rate in this model, since the interest rate is international). This shifts the $LM$ curve to the right, or at least tries to do so.
        \item fiscal stimulus of the government, i.e. the government increasing spending. This shifts the $IS$ curve to the right, or at least tries to do so.
        \item tariffs: the government can increase net exports through tariffs. This has the same effect as fiscal stimulus via an (attempted) right shift of the $IS$ curve.
        \item change of the world interest rate, which is now set on the global market or by the central bank of a large, dominant, reserve currency country abroad. This shifts the $ZZ$ curve upwards. 
    \end{itemize}
    }
    \showmytikz[1]{../shared/tikz/tikz-islmzzfindings}
\end{frame}

\begin{frame}{Fiscal Stimulus, Flexible Exchange Rate}
    \showmytikz[1]{../shared/tikz/tikz-mundellfleming}
    \mycomment{
    Fiscal stimulus shifts the $IS$-curve to the right. 
    
    Under flexible exchange rates this will result in an appreciation of the currency, and to capital import.
    
    The curve snaps right back as imports increase and exports fall. 
    
    The extra government spending has in effect lead to increased imports.
    
   As production shifts abroad, the policy results only in a total crowding-out of private activities.

   The finding is valid for tariffs. Tariffs do not work under flexible exchange rates.
    }
\end{frame}

\begin{frame}{Fiscal Stimulus, Fixed Exchange Rate}
    \showmytikz[1,2]{../shared/tikz/tikz-mundellfleming}
    \mycomment{

    If, however, the exchange rate is pegged, the central bank will sell home currency to keep the rate unchanged. This shifts the $LM$ curve to the right as well, resulting in higher output. 
    
    Thus, a fiscal stimulus can only work when the exchange rate is fixed. It works better than in the closed economy, as the interest rate remains stable.

    The finding is valid for tariffs. Tariffs work when the exchange rate is fixed.
    }
\end{frame}

\begin{frame}{Monetary Stimulus, Fixed Exchange Rate}
    \showmytikz[3]{../shared/tikz/tikz-mundellfleming}
    \mycomment{
    Monetary stimulus  shifts the LM curve to the right. The central bank can make this happen through open market operations only, since it cannot set the interest rate. It would be selling home currency.
    
    If the exchange rate is fixed, no such policy is possible. Any monetary stimulus will have to be offset by purchase of home currency by the central bank to hold the exchange rate stable. 
    
    Such a monetary policy contradicts itself, as the bank would have to expand and reduce the monetary base (buy and sell home currency) at the same time, which is impossible.

    This is known as Mundell's Trilemma. You cannot have free movement of capital, fixed exchange rates, and autonomous monetary policy. 
    }
\end{frame}

\begin{frame}{Monetary Stimulus, Flexible Exchange Rate}
    \showmytikz[3,4]{../shared/tikz/tikz-mundellfleming}
    \mycomment{
    Under flexible exchange rates, the open market expanse of the monetary base (selling of home currency) results in depreciation, leading to more exports and less imports.  As a result, the $IS$ curve also shifts to the right. The economy will have a larger output.

    Expansive open market policy only works when the exchange rate is flexible. It works better than in the closed economy, since the interest rate remains stable.
    }
\end{frame}


\begin{frame}{Change of Interest Rate}
    \showmytikz[1]{../shared/tikz/tikz-mundellfleming-zz}
    \mycomment{
    Finally, we look at the effect of changing international interest rates and the resulting imbalances. That is, $ZZ$ moves up. The new interest rate is too high for the local economy.
    }
\end{frame}


\begin{frame}{Change of Interest Rate, Fixed Exchange Rate}
    \showmytikz[1,2]{../shared/tikz/tikz-mundellfleming-zz}
    \mycomment{
    Under fixed exchange rates, the increased interest rate abroad has made internal investment less attractive. This means the currency comes under pressure as capital moves abroad. To keep the exchange rate, the central bank starts buying home currency. This shifts the $LM$-curve to the left, resulting in lower output at the higher interest rates. Investments will be reduced.
    }
\end{frame}

\begin{frame}{Change of Interest Rate, Flexible Exchange Rate}
    \showmytikz[1,3]{../shared/tikz/tikz-mundellfleming-zz}
    \mycomment{
    If the exchange rate is flexible, the rising international interest rate will have capital flow from inside the country to the rest of the world. This will depreciate the home currency. That results in less imports and more exports, thus higher output, that is a right shifting IS curve.

    The depreciation had a positive impact on the economy.

    There is a caveat. One result could be the import of inflation, since imported goods have now become more expensive relative to exported goods (terms of trade effect). The overall effect on the trade balance will depend on the elasticities of imports and exports. 
    }
\end{frame}

\begin{frame}{IS-LM-ZZ Findings: Overview}
    \mycomment{
    The results of the IS-LM-ZZ analysis show that
    \begin{itemize}
        \item under fixed exchange rates, only fiscal policy works. The hands of the central bank are tied; it must always act to stabilise the exchange rate. If there is fiscal stimulus, the central bank must support this by selling home currency to keep the peg.
        \item under flexible exchange rates, only monetary policy works. Any fiscal policy will be offset by changes in the exchange rate that lead to an offset through changes in the trade balance.
        \item rises of the international interest rate level will have a negative effect if the exchange rates are fixed, since the country must adapt that interest rate. If the exchange rate is flexible, a rise of the international exchange rate leads to depreciation and hence more output.
    \end{itemize}
    This can be formulated inversely:
    \begin{itemize}
        \item monetary stimulus requires flexible exchange rates,
        \item fiscal stimulus requires fixed exchange rates.
    \end{itemize}
    }
    \showmytikz[2]{../shared/tikz/tikz-islmzzfindings}
\end{frame}

\begin{frame}{IS-LM-ZZ and Tariffs}
    \mytext{
    Tariffs decrease imports, thus increase net exports, thus shift the IS curve to the right.

    That is just like a financial stimulus.

    Result: 
    
    Tariffs have no effect when the exchange rate is flexible. 

    Under fixed exchange rate, they have an effect, albeit to the detriment of the rest of the world, which might retaliate.
    }
    \mycomment{
        The short term effect of a tariff under flexible exchange rates is a rise of the (internal) interest rate, resulting in more capital import and then a rise of the exchange rate. This offsets the expansive effect perfectly in the medium run.

        However, there are distributional effects of the tariffs between the export and the import industry.

        Under fixed exchange rates, tariffs have a positive effect on the cost of the rest of the world (beggar-ty-neighbor policy). If they are countered by retaliation, everone loses out; this was the case in the Great Depression.
    }
\end{frame}

% \begin{frame}{Findings: Overview}
% \mytable{
%         \newcommand\exports{\text{Ex}}
%         \begin{tabular}{c||c|c|c}
%              Exch. Rate  & Monetary Policy & Fiscal Policy & Foreign Policy\\ 
%              Regime  & $\Delta i <0$ & $\Delta G >0$ & $\Delta i_{int} >0$ \\ 
%              & LM shifts right & IS shifts right  & ZZ shifts up\\
%              \hline\hline
%              % flexible & \textbf{affects GDP}  & primary effect & $\Delta Exp > 0$ \\
%             flexible     & depreciation & appreciation & depreciation \\  
%                  &  $\Delta \exports >0$ & $\Delta \exports <0$ & $\Delta \exports >0$\\
%                  & & offset & $IS$ shifts right\\
%                 & $\rightarrow$ effect & no effect & pos. effect\\
%              \hline
%              fixed & no depreciation  & no appreciation & no depreciation \\
%               & buy home curr. & sell home curr. & buy home curr. \\
%               & (contradiction) & $\rightarrow$ & $\rightarrow$ \\
%               &  & $LM$ shifts right & $LM$ shifts left\\
%               & & $\rightarrow$ effect & neg. effect \\
%         \end{tabular}
% }
% \end{frame}



\begin{frame}{Exchange Rate, Supply and Demand for Currency}
    \showmytikz{../shared/tikz/tikz-exchangerate}
    \mycomment{
    If my home currency (\euro) be infinitely strong ($\frac{\text{\euro}}{\$}=0$), I could get everything from abroad for free. But the value of all these imports in \euro\ would be zero (bottom left corner).

    As the currency depreciates, imports become more expensive, and I will import less and less. 
    
    At some point everything abroad will have become prohibitively expensive, I will stop importing, and my import value will again be zero (upper left corner).

    There is thus an exchange rate where the value of my imports in \euro\ is at maximum (Im$^{\text{\euro}}_{max}$, middle right). 
    }
\end{frame}

\begin{frame}{Exchange Rate, Supply and Demand for Currency}
    \showmytikz[2]{../shared/tikz/tikz-exchangerate}
    \mycomment{
        If my currency is super strong, I will not export anything (bottom left). As the currency devalues, I start to export. The cheaper my currency, the more I earn (in \euro ) as I export, and there is no end to this. 

        At some exchange rate the value of my exports will equal the value of my imports. At this point I have a balanced trade balance. This would be the market outcome in the long run, where we have purchasing power parity. 
        
        Every point above $E$ is an undervalued \euro, every point below an overvalued \euro.
    }
\end{frame}

\begin{frame}{Exchange Rate, Supply and Demand for Currency}
    \showmytikz[2,3]{../shared/tikz/tikz-exchangerate}
    \mycomment{
        If I manage to peg my currency at a rate lower than that equilibrium exchange rate, which in our graph is above the equilibrium, I can run a trade surplus. I will accumulate foreign currency. 

        There is an ideal trade surplus. If I depreciate further, I will not be getting more in return in terms of imports. My ideal trade surplus would be so that I realise Im$^{\text{\euro}}_{max}$.
        
        To keep the exchange rate at that point, the central bank must buy the foreign currency for local currency. This shifts the local LM curve. If that is not desired, the bank can sell bonds to reduce the money supply again (sterilisation). 
    }
\end{frame}

\begin{frame}{Sterilisation}
\mytext{
    \begin{itemize}
        \item Central Bank intervenes to keep peg
        \item Selling or buying own currency
        \item Selling own currency: depreciate
        \item Buying own currency: appreciate
        \item Buying is limited by available reserves
    \end{itemize}
}
    
\end{frame}

\begin{frame}{Marshall-Lerner-Condition}
    \mytext{
    
    Depreciation: good or bad for BoP?
    \begin{itemize}
        \item Depends on elasticities 
        \item of imports and export demand to exchange rate change
        \item hence on demand to exports and imports
        \item normally: depreciation good for BoP
    \end{itemize}
    }
\end{frame}
\begin{frame}{Marshall-Lerner-Condition}
    \mytext{
    For a depreciation to improve the trade balance (TB):
    \begin{equation}
        M = |\eta_{IM}| + |\eta_{Ex}| > 1
    \end{equation}

    \begin{description}
        \item[$M<1$] exchange market unstable 
        \item[$M=1$] depreciation no effect on TB
        \item[$M>1$] market stable, depreciation positive effect on TB
    \end{description}

    Note that elasticities must be pronounced for a depreciation to have any effect (other than inflation)
    }
\end{frame}

\begin{frame}{Financial Openness}
    \mytext{
    \begin{itemize}
        \item Chinn-Ito Index \cite{theGlobalEconomy.com}
        \item Wang-Jahan Index (0...1) (12 asset classes), IMF Capital Account Openness \cite{IMFopenness}
    \end{itemize}
    }
\end{frame}

\begin{frame}{Net International Investment Position (NIIP)}
    \showmybitmap[https://ec.europa.eu/eurostat/statistics-explained/index.php?title=\%20International\_investment\_position\_statistics]{../shared/bitmaps/Net_international_investment_position,_surplus_and_deficit,_Top_10_economies,_2024_(EUR_billion).png}
    \mycomment{
    \begin{itemize}
        \item NIIIP is a stock quantity
        \item often compared with GDP (but should be with capital stock)
        \item positive for most developed countries, \item negative for the US
    \end{itemize}
    }
\end{frame}

\begin{frame}{Net International Income}
    \mytext{
    \begin{itemize}
        \item NII is a flow quantity
        \item income from NIIIP (interest payments)
        \item positive for most developed countries, including the US
    \end{itemize}
    }   
\end{frame}

\begin{frame}{NIIP-NIP paradoxon}
    \begin{itemize}
        \item US is net debtor, but receives more income from debt than it pays
        \item China is net creditor, but pays more debt service than it receives
        \item Explanation I: `dark matter' i.e. stuff not accounted for: payments for immaterial assets?
        \item Explanation II: riskier and not so risky assets: interest rate spread
        \item US invests in risky, high-yield assets
        \item rest of the world invests in T-bills (US treasury bonds) with little yield
    \end{itemize}
    Read: \cite{schmitt-grohe-slides}
\end{frame}

%%%%%%%%%%%%%%%%%%%%%%%%%%%%%%%%%%%%%%%

